% !TEX root = ../../main.tex
\documentclass[../../main.tex]{subfiles}

\begin{document}

\chapter{Unpacking The Socio-Materiality of Datacentering for Sustainability Transitions}
\labch{datacentering}

\begin{fullwidthpar}
\textit{This perspective article, prepared for Environmental Innovation and Societal Transitions, introduces ``datacentering'' as a concept for studying digital infrastructure through sustainability transitions frameworks. It proposes multimodal cartography as a methodology for circumventing the opacity of the data center industry---an approach that directly informs the methodological innovations developed in this dissertation.}
\end{fullwidthpar}

\section{Introduction: The Materiality Behind the Cloud}

In 2024, global data center electricity consumption surpassed 400 TWh---more than the entire national consumption of France \sidecite{IEA2024}. Microsoft's water usage increased 34\% in a single year, driven largely by data center cooling \sidecite{Mytton2021}. Training a single large language model can emit as much carbon as five automobiles over their lifetimes \sidecite{Strubell2019}. As artificial intelligence reshapes industries worldwide, these figures are projected to grow dramatically: the International Energy Agency estimates data center energy demand could double by 2030.

These are not marginal externalities. They represent the material foundations of what we casually call ``the cloud.''

\begin{kaobox}[title=The Cloud Imaginary]
The cloud metaphor has proven remarkably durable---and remarkably misleading. It functions as what Jasanoff and Kim term a \textit{sociotechnical imaginary}: a collectively held vision of desirable futures that shapes technological development and legitimates particular arrangements of power. The cloud imaginary conjures images of weightless, placeless computation: data flowing through an immaterial digital ether.
\end{kaobox}

Yet this framing systematically obscures one of the most significant infrastructural transformations of our time. Data centers are massive physical installations that consume extraordinary quantities of energy, water, and land while reshaping territories, labor markets, and ecologies worldwide \sidecite{Monstadt2025, Carruth2014}.

Despite mounting evidence of these impacts, sustainability transitions research has largely treated digital infrastructure as an \textit{enabler} of other transitions---facilitating smart grids, precision agriculture, and circular economy platforms---rather than as a socio-technical system requiring transformation in its own right \sidecite{Vial2019}. This perspective argues for a fundamental reorientation.

We introduce the concept of \textbf{datacentering} to denote the active processes through which digital infrastructure is materially emplaced, governed, and contested. Datacentering is not merely about data centers as discrete facilities. It encompasses the relational practices, politics, and power dynamics through which digital infrastructure reshapes territories and communities.\marginnote{The term shifts attention from static objects to dynamic processes---from ``data centers'' as nouns to ``datacentering'' as ongoing socio-material transformation.}

Advancing this agenda requires methodological innovation. The data center industry is characterized by corporate secrecy, non-disclosure agreements, and deliberate invisibility. Conventional research methods---interviews with executives, facility tours, official documentation---face significant access barriers. We therefore propose a \textit{multimodal cartographic approach} that triangulates diverse evidence sources: corporate disclosures, social media platforms, news archives, satellite imagery, regulatory filings, and visual media.

%----------------------------------------------------------
\section{Bibliometric Landscape: A Fragmented Field}
%----------------------------------------------------------

Scholarship on data centers is growing rapidly but remains deeply fragmented. To map this intellectual terrain and identify opportunities for transitions research, we conducted a bibliometric analysis of Scopus-indexed literature from 2010 to 2025.

\subsection{Findings}

\textbf{Rapid growth, recent acceleration.} Publication volumes have grown exponentially since 2015, with marked acceleration after 2018---coinciding with heightened public awareness of AI's environmental footprint and journalistic investigations into Big Tech's resource consumption \sidecite{Crawford2021}.

\textbf{Disciplinary silos.} The field is dominated by engineering and computer science, focused primarily on technical efficiency: Power Usage Effectiveness (PUE) optimization, cooling innovations, and energy management. A smaller but growing body of critical scholarship emerges from geography, sociology, science and technology studies (STS), and media studies, examining data centers as sites of socio-political contestation. These two literatures rarely cite each other.

\textbf{Three disconnected clusters.} Keyword co-occurrence analysis reveals three distinct research communities with limited interconnection:

\begin{itemize}
    \item \textit{Technical efficiency cluster}: PUE, cooling, energy efficiency, optimization, workload management
    \item \textit{Environmental impact cluster}: carbon footprint, water usage, lifecycle assessment, renewable energy, e-waste
    \item \textit{Socio-political cluster}: governance, territory, justice, community, labor---the smallest and least integrated cluster
\end{itemize}

\subsection{Critical Gaps}

Our analysis reveals three gaps that motivate the research agenda we propose:

\begin{enumerate}
    \item \textbf{Absent transitions frameworks.} The Multi-Level Perspective, Technological Innovation Systems, and Strategic Niche Management have rarely been applied to digital infrastructure \sidecite{Geels2002, Markard2012}.
    \item \textbf{Underdeveloped spatial analysis.} Despite cartography's proven utility in other infrastructure domains, spatial and geographic approaches to datacentering remain nascent.
    \item \textbf{Methodological barriers.} Industry opacity---corporate secrecy, restricted access, NDA cultures---constrains conventional research methods.
\end{enumerate}

%----------------------------------------------------------
\section{Conceptual Framing: Three Streams, One Synthesis}
%----------------------------------------------------------

Our research agenda synthesizes three theoretical streams that have developed largely in isolation: critical data center studies, socio-materiality theory, and sustainability transitions scholarship.

\begin{table}[h]
\caption{Theoretical synthesis: Three streams informing the research agenda}
\labtab{theoretical}
\begin{tabular}{p{2.5cm}p{3cm}p{2.5cm}p{2.5cm}}
\toprule
\textbf{Stream} & \textbf{Core insight} & \textbf{Key concepts} & \textbf{Contribution} \\
\midrule
Critical Data Center Studies & Data centers are political sites & Territorial politics, timescapes & Substantive knowledge \\
\addlinespace
Socio-Materiality & Digital and material are inseparable & Entanglement, data journeys & Analytical lens \\
\addlinespace
Sustainability Transitions & Systemic change requires multi-level analysis & Regimes, niches, landscapes & Transformation pathways \\
\bottomrule
\end{tabular}
\end{table}

\subsection{Critical Data Center Studies}

Over the past decade, scholars across media studies, geography, anthropology, and STS have established what Edwards et al.\ term ``Critical Data Center Studies'' \sidecite{Edwards2025}. This interdisciplinary field treats data centers not as neutral technical facilities but as sites where environmental, social, economic, and political forces converge.

Three themes organize this literature:

\textbf{Territorial politics.} Data center siting involves complex negotiations among corporations, states, and communities. Operators seek cheap electricity, fiber connectivity, tax incentives, and political stability; localities seek jobs and investment but may resist environmental burdens.

\textbf{Temporal dimensions.} Velkova and Plantin introduce ``timescapes'' to analyze how data centers operate across multiple temporalities: the milliseconds of computation, the years of facility lifecycles, the decades of environmental slow violence from resource extraction and e-waste accumulation \sidecite{Velkova2023}.

\textbf{Metabolic relations.} Data centers are metabolic systems, consuming energy, water, and land while producing heat, carbon emissions, and electronic waste.

\subsection{Socio-Materiality}

Socio-materiality, as developed in organization studies and STS, theorizes the constitutive entanglement of social and material dimensions \sidecite{Orlikowski2007}. Applied to digital infrastructure, this perspective insists that ``the cloud'' cannot be understood apart from its material substrates: the cables, chips, cooling systems, buildings, and human bodies that make computation possible.

\subsection{Sustainability Transitions Theory}

We propose that the data center sector constitutes a \textbf{regime in formation}. It exhibits classic regime characteristics---technological lock-in, incumbent power concentration, path dependency---while remaining more fluid than established infrastructure domains like energy or transport. Three hyperscalers (Amazon, Microsoft, Google) dominate the market. Proprietary architectures create switching costs. Massive sunk investments in facilities and fiber create path dependencies.

\begin{kaobox}[title=The Dual Role of Data Centers]
Data centers occupy a paradoxical position: they are simultaneously infrastructure \textit{for} sustainability transitions (enabling smart grids, precision agriculture, dematerialization) and infrastructure \textit{requiring} transition (toward decarbonization, circularity, and environmental justice).
\end{kaobox}

%----------------------------------------------------------
\section{Toward a Multimodal Cartographic Research Agenda}
%----------------------------------------------------------

We propose \textbf{multimodal cartography} as a methodological framework for studying the socio-materiality of datacentering. This approach combines the spatial sensibilities of critical cartography with the evidentiary breadth of multimodal research design.

\subsection{Why Cartography?}

Mapping offers four distinct affordances for infrastructure research:

\begin{margintable}
\caption{Cartographic affordances}
\begin{tabular}{p{2cm}p{3cm}}
\toprule
\textbf{Affordance} & \textbf{Application} \\
\midrule
Rendering visible & Countering cloud rhetoric \\
\addlinespace
Relational geography & Tracing connections \\
\addlinespace
Counter-practice & Democratic contestation \\
\addlinespace
Temporal dynamics & Documenting change \\
\bottomrule
\end{tabular}
\end{margintable}

Data centers are designed for invisibility: windowless buildings, remote locations, generic architecture, and cloud metaphors all work to dematerialize digital infrastructure in public consciousness. Critical cartography counters this by revealing spatial patterns of resource consumption, environmental impact, and territorial transformation.

\subsection{Why Multimodal?}

Traditional infrastructure research relies on interviews, ethnography, and official documents. These methods face significant barriers when studying an industry characterized by corporate secrecy, NDA-bound employees, restricted facility access, and strategic opacity.

Multimodal cartography addresses these challenges by treating diverse media as spatial data sources:

\begin{itemize}
    \item \textbf{Corporate digital traces}: Websites, sustainability reports, investor presentations, SEC filings
    \item \textbf{Social media}: Twitter/X, LinkedIn, Reddit, community groups
    \item \textbf{News archives}: Local press, national media, trade publications
    \item \textbf{Visual/satellite media}: Satellite imagery, YouTube, photographs
    \item \textbf{Regulatory documents}: EIAs, permits, utility filings, tax agreements
    \item \textbf{Investment data}: Crunchbase, PitchBook, corporate registries
\end{itemize}

\subsection{The Logic of Triangulation}

The power of multimodal cartography lies in \textbf{triangulation}---systematically cross-referencing evidence across sources to build robust accounts of datacentering dynamics.\marginnote{This triangulation approach directly parallels the multimodal methodology developed in \RoleBox for mapping role constellations across different data modalities.}

Consider a hypothetical example: A corporate sustainability report claims a new facility will achieve a Power Usage Effectiveness (PUE) of 1.2 and use 100\% renewable energy. Satellite imagery reveals the facility is connected to a coal-heavy regional grid. Local news archives document community opposition to water consumption projections. Social media posts from facility workers describe diesel backup generators running frequently. Regulatory filings show the operator received a 20-year tax abatement worth \$50 million.

No single source tells the full story.

%----------------------------------------------------------
\section{Research Directions: Seven Cartographic Domains}
%----------------------------------------------------------

We outline seven research directions that operationalize multimodal cartography for datacentering studies:

\subsection{Metabolic Cartographies: Energy, Water, and Land}

Data centers are metabolic systems with significant environmental footprints. A single hyperscale facility can consume as much electricity as a small city and millions of gallons of water daily for cooling.

\textit{Real-world context:} In 2022, residents of The Dalles, Oregon discovered that Google's data centers consumed over a quarter of the city's water supply---information that emerged only through public records requests after years of corporate opacity.

\subsection{Governance Geographies: Regulation, Incentives, and Policy Mobility}

Data center siting unfolds within fragmented governance landscapes where jurisdictions compete through tax incentives, relaxed regulations, and streamlined permitting.

\textit{Real-world context:} Virginia's Loudoun County, hosting the world's largest concentration of data centers, has foregone over \$1 billion in tax revenue through incentive programs. Meanwhile, Amsterdam and Singapore have implemented moratoria on new data center construction.

\subsection{Supply Chain Cartographies: Extraction, Production, Disposal}

Data centers are nodes in extended material networks spanning extraction frontiers, manufacturing hubs, and e-waste destinations. A server's lifecycle connects rare earth mining in the Democratic Republic of Congo, chip fabrication in Taiwan, assembly in China, operation in Virginia, and eventual disposal in Ghana or Pakistan.

\subsection{Contestation Cartographies: Discourse, Resistance, and Power}

Datacentering is not merely technical but deeply political. Communities contest facility siting; activists challenge corporate sustainability claims; workers organize around labor conditions.

\textit{Real-world context:} In 2023, residents of Beaverton, Oregon successfully blocked a proposed Meta data center, citing concerns about noise, water consumption, and neighborhood character.

\subsection{Labor and Community Cartographies: Work, Place, and Justice}

Data centers promise economic development but often deliver fewer jobs than anticipated---highly automated facilities may employ only 50--100 workers despite multi-billion-dollar investments.

\subsection{Investment and Corporate Cartographies: Capital, Power, and Ownership}

The data center sector has attracted massive capital inflows, with private equity, sovereign wealth funds, and pension funds joining hyperscalers in financing infrastructure expansion.

\textit{Real-world context:} Blackstone, the world's largest private equity firm, has invested over \$50 billion in data center infrastructure globally.

\subsection{Transition Pathway Cartographies: Alternatives and Futures}

Beyond critique, multimodal cartography can support envisioning and planning alternative datacentering futures. Experiments with community-owned infrastructure, circular economy models, and demand-side governance offer glimpses of different possibilities.

\textit{Real-world context:} Barcelona's ``technological sovereignty'' initiative explores municipal and cooperative alternatives to hyperscaler dominance.

%----------------------------------------------------------
\section{Discussion: Implications for Research, Policy, and Practice}
%----------------------------------------------------------

\subsection{Theoretical Contributions}

\textbf{Digital infrastructure as a transitions domain.} Our first and most fundamental argument is that digital infrastructure deserves the same analytical attention from transitions scholars as energy, transport, agriculture, or housing. The data center sector exhibits classic regime characteristics: technological lock-in through proprietary architectures and switching costs; incumbent power concentrated in three hyperscalers; institutional embedding via tax incentives and regulatory capture; and path dependency from massive sunk investments.

\textbf{Spatial sensitivity in transitions research.} Sustainability transitions scholarship has been critiqued for insufficient attention to space and geography. Datacentering offers a productive site for developing more spatially attentive frameworks.

\textbf{Methodological innovation for opaque sectors.} The data center industry's secrecy is not unique. Similar challenges confront researchers studying platform companies, financial infrastructure, logistics networks, and other ``hidden'' systems central to contemporary capitalism. Multimodal cartography offers a transferable toolkit for studying sectors characterized by corporate opacity, restricted access, and strategic invisibility.

\subsection{Policy Implications}

Our analysis suggests several directions for policymakers:

\begin{enumerate}
    \item \textbf{Mandatory disclosure requirements.} Jurisdictions should require comprehensive, standardized reporting of data center energy consumption, water usage, and carbon emissions.
    \item \textbf{Conditional incentives.} Tax abatements should be tied to enforceable sustainability commitments.
    \item \textbf{Cumulative impact assessment.} Environmental review should assess cumulative impacts of data center clusters.
    \item \textbf{Participatory governance.} Affected communities should have meaningful voice in siting decisions.
    \item \textbf{International coordination.} Given the global mobility of data center capital, coordination is needed to prevent regulatory arbitrage.
\end{enumerate}

\subsection{Conclusion}

The proliferation of data centers represents one of the most consequential infrastructural transformations of our era---yet one that remains remarkably understudied from a sustainability transitions perspective. The ``cloud'' metaphor has done ideological work in obscuring this transformation, but the material realities of datacentering are increasingly impossible to ignore.

We have proposed \textbf{datacentering} as a concept that foregrounds these active socio-material processes, and \textbf{multimodal cartography} as a methodological framework for their study. As cities from Amsterdam to Dublin to Singapore impose moratoria on new facilities, as communities from Oregon to Ireland resist data center siting, and as the environmental footprint of AI becomes undeniable, there is urgent need for research that can inform more just and sustainable digital futures.

\begin{kaobox}[title=The Cloud Has a Geography]
The infrastructural foundations of our digital world are being built now---in specific places, with specific environmental and social consequences, benefiting specific actors. Mapping datacentering, in all its metabolic, governmental, contested, and speculative dimensions, is a necessary step toward ensuring that these foundations do not foreclose the possibility of sustainability transitions writ large.
\end{kaobox}

\end{document}
